\textbf{Статус.} Гипотеза~\ref{conj:master} воспроизводит:
(1) все 23 вычисленные точки $P_n(q)$ из факта~\ref{fact:table};
(2) все по-конфигурационные пересчёты (свыше $4800$ пар <<конфигурация
$\mapsto$ число идеалов>> при $n\leq 9$);
(3) решающий тест: для конфигурации
$L^{*}=\{(2,1),(4,2),(5,3),(7,5),(8,6),(9,8)\}$ в $\NT(9,2)$ (две
группы склеек по две связанные пары) конкурирующая модель без слияний групп
предсказывала $1646$ идеалов, гипотеза~\ref{conj:master} "--- $1672$;
два независимых прямых перечисления (обход накрытий через цоколь фактора и
редукция форсированным подпространством) дали $1672$.
Из гипотезы следуют явные многочлены

%% ==== вырезано при замене conj:master теоремой ====

\item Доказательство гипотезы~\ref{conj:master}: полнота классификации
склеек (частные случаи "--- вклад одной связанной пары $(q-1)W_m$ и
пример \S\,\ref{sec:nt4} "--- допускают прямую проверку).
\item Полиномиальность $P_n(q)$ по $q$.
